\documentclass[11pt]{article}

\usepackage[margin=1in]{geometry}
\usepackage[T1]{fontenc}
\usepackage{lmodern}
\usepackage{amsmath,amssymb}
\usepackage{booktabs}
\usepackage{graphicx}
\usepackage{tabularx}
\usepackage{longtable}
\usepackage{array}
\usepackage[table]{xcolor}
\usepackage{hyperref}
\usepackage{microtype}
\usepackage{caption}
\usepackage{subcaption}
\usepackage{enumitem}
\usepackage{fancyhdr}

\graphicspath{{cognitive_gaze_figures/}{./}}
\hypersetup{
  colorlinks=true,
  linkcolor=blue!45!black,
  citecolor=blue!45!black,
  urlcolor=blue!55!black
}

\definecolor{Navy}{HTML}{122B3E}
\definecolor{Teal}{HTML}{007D84}
\definecolor{Blue}{HTML}{316FAA}
\definecolor{Amber}{HTML}{BC7918}
\definecolor{Red}{HTML}{AD413D}
\definecolor{Light}{HTML}{D3D8DA}
\definecolor{PaleTeal}{HTML}{E2F2F1}
\definecolor{PaleBlue}{HTML}{E7EFF7}
\definecolor{PaleAmber}{HTML}{FBF1D9}
\definecolor{PaleRed}{HTML}{F5DEDC}

\setlength{\parindent}{0pt}
\setlength{\parskip}{0.65em}
\setlist[itemize]{topsep=0.25em,itemsep=0.2em}
\setlist[enumerate]{topsep=0.25em,itemsep=0.2em}
\pagestyle{fancy}
\fancyhf{}
\lhead{\small SD-SLO gaze tracking from first principles}
\rhead{\small \thepage}
\renewcommand{\headrulewidth}{0.4pt}

\newcommand{\evidence}[1]{\textsc{\color{Teal}#1}}
\newcommand{\missing}{\textsc{\color{Red}missing}}
\newcommand{\hypothesis}{\textsc{\color{Red}hypothesis}}
\newcommand{\fixture}{\textsc{\color{Amber}fixture}}
\newcommand{\measured}{\textsc{\color{Teal}measured}}
\newcommand{\derived}{\textsc{\color{Blue}derived}}
\newcommand{\programfact}{\textsc{\color{Navy}program fact}}

\title{\textbf{From Scanned Photons to Time-Resolved Retinal Gaze}\\
\large A First-Principles Engineering and Evidence Audit of an SD-SLO Cognitive Platform}
\author{Kodiak Sciences gaze-tracking program}
\date{22 July 2026}

\begin{document}
\maketitle

\begin{abstract}
This document explains how a scanning laser ophthalmoscope (SLO), originally optimized for
retinal imaging, can be used to estimate eye motion and support synchronized cognitive experiments.
It is written for a technically literate reader who does not already know the instrument.
The analysis begins with the optical and electronic hardware, derives the scanner and sampling
clocks, and then follows one retinal displacement estimate through rectification, strip formation,
preprocessing, normalized cross-correlation, quality gating, subpixel placement, calibration,
and trajectory assembly. Numerical statements are separated into measured values, explicit
derivations, literature reports, synthetic fixtures, program facts, hypotheses, and missing inputs.

For the candidate 10 kHz sinusoidal fast axis, bidirectional use produces a 20 kline/s clock and a
50 $\mu$s one-way sweep. A 24-line strip therefore spans 1.20 ms; a 6-line stride places correlated
outputs every 0.30 ms during active acquisition. Existing synthetic experiments demonstrate
sub-arcminute registration under their declared model, while real retinal data demonstrate
lock behavior and coarse pursuit correspondence without providing independent high-rate eye-motion
truth. Two conditional timing models expose intervals not explained by active lines, but they must
not be combined: the legacy capture assumes 58.34 ms active plus 10 ms dead time, while applying the
candidate 20 kline/s clock to 808 legacy lines yields 40.4 ms active plus 27.94 ms unexplained at
14.633 fps. Neither identifies physical scanner flyback without telemetry.

The paper also separates biological event rate, waveform bandwidth, acquisition rate, and spatial
precision for saccades, microsaccades, drift, tremor, and pupil dynamics. Tremor is intrinsically
arcsecond-scale, but high-rate drift increments and fine direction estimates for the smallest
microsaccades can also become arcsecond-limited. Published adaptive-optics SLO performance cannot be
transferred to this conventional SD-SLO without validation. Continuous optical operation is treated
through a parameterized exposure framework only: laser-product classification, ophthalmic patient
exposure, and medical-device classification are distinct evaluations, and the measured source and
fault inputs required for a safety conclusion are presently incomplete.
\end{abstract}

\textbf{Keywords:} scanning laser ophthalmoscopy; retinal image registration; eye tracking;
normalized cross-correlation; fixational eye motion; scanner timing; laser safety; OPM-MEG;
cognitive instrumentation.

\tableofcontents
\newpage

\section{Purpose, scope, and evidence discipline}

\subsection{Why the system requires a new architecture review}

The program began with a retinal-health objective: a portable SLO or OCT device that could acquire
high-quality retinal images outside the clinic and provide denser longitudinal information than
monthly visits. That objective naturally prioritizes image quality, patient workflow, portability,
and intermittent acquisition. A frame can be judged successful if it is clinically interpretable,
even when time between frames is not observed.

The emerging objective is different. The device is now being considered as a cognitive and
attention measurement platform that can synchronize retinal motion, pupil diameter, visual
stimulus timing, behavior, and OPM-MEG. OPM coexistence is accepted here as a characterized
\programfact; the production requirement is to bind that result to controlled hardware,
firmware, and experimental configurations. Cognitive experiments place much greater weight on
event time, continuity, uncertainty, and explicit missingness. A 10 ms unobserved interval can be
negligible for a still retinal image but decisive for microsaccade onset or direction.

The architecture question is therefore not whether the prior medical design was ``wrong.'' Many
of its features---confocal retinal imaging, optical precision, calibration discipline, safety
controls, and portable mechanics---remain valuable. The question is which assumptions were
appropriate only for intermittent imaging and should now be changed, made configurable, or held
open before production hardware is frozen.

\subsection{Evidence classes}

Every quantitative or factual claim in this paper is assigned one of the following roles.

\begin{table}[ht]
\centering
\caption{Evidence classes used throughout the analysis.}
\begin{tabularx}{\textwidth}{>{\bfseries}p{0.17\textwidth} X}
\toprule
Class & Meaning\\
\midrule
\measured & Directly measured on the stated hardware, file, or capture.\\
\derived & Calculated from cited inputs using an explicit formula.\\
Literature & Reported by an external primary or authoritative source.\\
\fixture & Known-reference software or synthetic experiment; not biological validation.\\
\programfact & Project history or capability supplied by the program and accepted as true.\\
\hypothesis & Candidate design value or scientific interpretation requiring validation.\\
\missing & Required value is absent; no numerical conclusion is permitted.\\
\bottomrule
\end{tabularx}
\end{table}

This policy prevents four common category errors: synthetic known truth is not silently promoted to
human validation; predicted samples are not promoted to observations; nominal eye geometry is not
promoted to subject-specific calibration; and a symbolic exposure equation is not promoted to a
safety determination. The machine-readable source of truth is
\path{cognitive_gaze_evidence.json}. Generated figures are bound to that ledger, the generator,
and input hashes in \path{cognitive_gaze_figure_manifest.json}.

\section{Instrument primer: how an SLO creates retinal data}

\subsection{Optical and electronic path}

A scanning laser ophthalmoscope is a point-scanning retinal camera. An illumination source is
focused to a small retinal spot. A fast scanner moves that spot across one image axis, while a slow
scanner advances the line position along the other axis. Returning retinal light is descanned,
detected, amplified, and digitized by an analog-to-digital converter (ADC). The fundamental
measurement is therefore not a two-dimensional image but a time series:
\[
  \mathcal D = \{I(t_k), x_s(t_k), y_s(t_k), b(t_k), r_k\}_{k=1}^{N},
\]
where $I$ is detector intensity, $(x_s,y_s)$ is scanner position, $b$ is illumination/blanking
state, and $r_k$ collects validity and revision metadata. Image formation maps those samples to a
retinal coordinate grid.

The distinction matters because eye motion occurs while the image is being scanned. If $R(x,y)$ is
a stationary retinal reflectance pattern and $\mathbf d(t)$ is eye displacement, the detector can
be approximated as
\[
  I(t) \approx R\!\left(x_s(t)-d_x(t),\,y_s(t)-d_y(t)\right)+\eta(t),
\]
with noise and unmodeled contrast represented by $\eta(t)$. Early and late parts of one nominal
frame therefore correspond to different eye positions. Motion is written into the raster as vessel
shear and displacement.

\subsection{Fast and slow axes}

The candidate fast scanner follows a 10 kHz full sinusoid. One complete cycle is
\[
  T_{\mathrm{sine}}=\frac{1}{10{,}000}=100~\mu\mathrm{s}.
\]
If both directions are acquired, each half-cycle is a line:
\[
  T_{\mathrm{line}}=50~\mu\mathrm{s},\qquad
  f_{\mathrm{line}}=20{,}000~\mathrm{lines/s}.
\]
The fast axis does not move uniformly. With retinal amplitude $A$,
\[
  x(t)=A\sin(2\pi f t+\phi),\qquad
  v(t)=2\pi fA\cos(2\pi f t+\phi).
\]
Speed is maximal at field center and approaches zero at the turnarounds. Uniform-time ADC samples
are consequently sparse in retinal distance near field center and dense near the edges.

The slow scanner advances from line to line. A frame is a software grouping of these sequential
measurements, not an instantaneous exposure. For 808 lines at the candidate bidirectional rate,
\[
  T_{808} = \frac{808}{20{,}000}=40.4~\mathrm{ms}.
\]
This is a \derived{} active-line time. It does not reveal the present slow-axis return waveform,
blanking, or optical duty cycle.

\begin{figure}[ht]
\centering
\includegraphics[width=\textwidth]{01_scan_physics.png}
\caption{Fast-axis kinematics and sampling. The ADC samples uniformly in time, while scanner
velocity makes sample spacing nonuniform in retinal distance.}
\label{fig:scan}
\end{figure}

\subsection{Retinal scale is a model until it is calibrated}

The technical note uses a 24 mm nominal axial length and a 48$^\circ$ optical field. The Bennett
retinal magnification model \cite{Bennett1994} gives
\[
  q=0.01306(AL-1.82)=0.2896708~\mathrm{mm/deg},
\]
and hence a nominal field width
\[
  W=48q=13.9042~\mathrm{mm}.
\]
The arithmetic is \derived, but the 24 mm eye and stated field are \hypothesis{} design inputs.
They must not be interpreted as measured subject anatomy. Image displacement can be estimated in
pixels without this model; conversion to degrees, arcminutes, or millimeters requires explicit
spatial calibration and, for physical retinal distance, eye-specific biometry.

\section{Bandwidth, sampling, and spatial information}

\subsection{Analog bandwidth becomes spatial blur}

The measured end-to-end analog bandwidth at the ADC input is 10 MHz and is modeled in the technical
note as a one-pole response. The time constant and 10--90\% rise time are
\[
  \tau=\frac{1}{2\pi B}=15.915~\mathrm{ns},\qquad
  t_{10\text{--}90}=\ln(9)\tau=34.970~\mathrm{ns}.
\]
At the modeled center velocity of 436.8 m/s, this corresponds to a 15.275 $\mu$m electronic edge.
Convolving the electronic response with the assumed optical line-spread function yields an
18.034 $\mu$m combined center edge in the note's model. The bandwidth is measured; the one-pole
shape and 10 $\mu$m Airy first-zero diameter are modeling assumptions.

The key point is that ``more stored pixels'' cannot restore analog information. A denser output
grid can make interpolation smoother, but optical blur, detector bandwidth, ADC rate, and noise set
the recoverable information.

\subsection{ADC samples and output grids}

A 50 MS/s candidate ADC acquires
\[
  N_{\mathrm{raw}}=(50\times10^6)(50\times10^{-6})=2500
\]
samples per one-way sweep. At field center the modeled raw spacing is
\[
  \Delta x_{\mathrm{center}}=\frac{436.8~\mathrm{m/s}}{50\times10^6~\mathrm{s^{-1}}}
  =8.736~\mu\mathrm{m}.
\]
Output-grid pitch should therefore be interpreted relative to both raw sampling and the modeled
18.034 $\mu$m edge. A 1600-column image may be useful for interpolation and registration but does
not imply 1600 independent optical resolution elements.

\begin{figure}[ht]
\centering
\includegraphics[width=\textwidth]{02_resolution_sampling.png}
\caption{Bandwidth, noise, sampling, and output-grid trade-offs. Increasing bandwidth sharpens the
modeled edge but increases integrated noise; increasing output columns reduces grid pitch but does
not create new optical information.}
\label{fig:resolution}
\end{figure}

\subsection{Rectification or ``desinusoiding''}

Desinusoiding maps uniform-time samples to an ascending retinal coordinate. If an ideal sinusoid is
used, the inverse coordinate map is derived from $x(t)$. In production, a measured position
look-up table (LUT) should replace the ideal model. Detector delay must be corrected before
interpolation, because a time offset becomes a direction-dependent spatial error. Forward and
reverse lines must be mapped to one declared coordinate orientation, while turnaround, blanked, and
invalid samples remain explicit.

Rectification changes coordinates; it cannot recover photons, bandwidth, or samples that were not
acquired. This is an important boundary when discussing frame gaps or scanner return.

\section{From retinal texture to a gaze trajectory}

\subsection{Why retinal registration can measure eye motion}

Retinal vessels and other stable structures provide a texture coordinate system. Over a sufficiently
short interval, a strip extracted from the moving raster should resemble a translated region of a
reference retina. Estimating the translation that maximizes similarity produces one displacement
sample. Repeating the operation over overlapping strips creates a high-rate relative trajectory.

This measurement model assumes that local appearance is sufficiently stable after preprocessing,
that motion within a strip is not too large, that the reference contains the corresponding field,
and that the correlation peak is identifiable. Violations should produce explicit rejected states,
not silently filled positions.

\subsection{Strip aperture and output spacing}

Strip width and stride control different physical quantities. For width $W_s$, stride $S_s$, and
line rate $f_L$,
\[
  T_{\mathrm{aperture}}=\frac{W_s}{f_L},\qquad
  T_{\mathrm{spacing}}=\frac{S_s}{f_L},\qquad
  f_{\mathrm{output}}=\frac{f_L}{S_s}.
\]
For the 24/6 example at 20 kline/s,
\[
  T_{\mathrm{aperture}}=1.20~\mathrm{ms},\qquad
  T_{\mathrm{spacing}}=0.30~\mathrm{ms},\qquad
  f_{\mathrm{output}}=3.33~\mathrm{kHz}.
\]
Adjacent estimates share $1-S_s/W_s=75\%$ of their lines. Thus 3.33 kHz is the active output
spacing, not an independent photon-information rate. Wider strips usually improve registration
signal-to-noise at the cost of temporal averaging; smaller stride increases localization density at
the cost of stronger correlation among outputs.

\begin{figure}[ht]
\centering
\includegraphics[width=\textwidth]{04_strip_overlap.png}
\caption{Overlapping strips. Width sets the temporal aperture; stride sets the spacing between
successive, correlated estimates.}
\label{fig:strips}
\end{figure}

\subsection{Preprocessing and normalized cross-correlation}

Broad illumination variation can dominate raw intensity without helping registration. The current
workflow uses operations such as debanding, local-mean removal, and difference-of-Gaussians (DoG)
filtering to emphasize vessel-scale structure. These operations improve the numerical matching
problem; they do not improve the biological optical resolution.

For template $T$ and reference search region $R$, normalized cross-correlation (NCC) at displacement
$(u,v)$ is
\[
  \operatorname{NCC}(u,v)=
  \frac{\sum_{x,y}[T(x,y)-\bar T][R(x+u,y+v)-\bar R_{u,v}]}
  {\sqrt{\sum_{x,y}[T(x,y)-\bar T]^2
  \sum_{x,y}[R(x+u,y+v)-\bar R_{u,v}]^2}}.
\]
Local mean subtraction suppresses additive brightness offsets, while energy normalization makes
the score less sensitive to local contrast magnitude. The response over all $(u,v)$ is a 2-D
correlation surface. Peak location gives displacement, peak height is one quality measure, peak
shape exposes ambiguity, and a boundary peak can indicate that motion exceeded the search region.

The implementation refines the integer response maximum independently in $x$ and $y$ using a
three-point parabolic vertex. This estimates the local maximum of the sampled response; it does not
create optical resolution. A displacement is accepted only when acquisition validity and
registration quality pass their respective gates.

\begin{figure}[ht]
\centering
\includegraphics[width=\textwidth]{03_pipeline_walkthrough.png}
\caption{End-to-end strip registration: retinal data, strip selection, preprocessing, bounded
search, NCC surface, and time-stamped displacement.}
\label{fig:pipeline}
\end{figure}

\begin{figure}[ht]
\centering
\includegraphics[width=\textwidth]{05_ncc_placement.png}
\caption{One accepted correlation peak places a moving vessel strip into the reference coordinate
system. The displayed integer peak is for visual clarity; production uses subpixel refinement.}
\label{fig:placement}
\end{figure}

\subsection{Time assignment, calibration, and coordinate frames}

One trajectory row should contain more than a pair of coordinates. A minimal representation is
\[
  z_i =
  \left[t_i,\Delta x_i,\Delta y_i,q_i,a_i,e_i,p_i,c_i,r_i\right],
\]
where $t_i$ is the strip-center hardware time, $q_i$ is registration quality, $a_i$ is acquisition
state, $e_i$ is estimate state, $p_i$ records prediction provenance, $c_i$ identifies calibration,
and $r_i$ identifies the hardware/software revision. Strip-center timing acknowledges that each
estimate integrates over a finite aperture.

The relative pixel displacement must then pass through an explicit transform chain:
\[
  \Delta \mathbf p_{\mathrm{raw}}
  \xrightarrow{\text{direction/LUT}}
  \Delta \mathbf p_{\mathrm{rect}}
  \xrightarrow{\mathbf C_{\mathrm{tracker}}}
  \Delta \mathbf p_{\mathrm{retina}}
  \xrightarrow{\mathbf C_{\mathrm{spatial}}}
  \Delta \boldsymbol\theta_{\mathrm{eye}}.
\]
The chain includes forward/reverse line orientation, scanner-position LUT, detector delay,
cross-axis coupling, pixel anisotropy, sign convention, and spatial scale. Reference recentering or
replacement must be logged because it changes the coordinate origin.

\subsection{Observed, rejected, predicted, and missing are not interchangeable}

The production data model should preserve at least four estimate states:

\begin{itemize}
  \item \textbf{OBSERVED}: photons were acquired and registration passed acceptance criteria;
  \item \textbf{REJECTED}: photons were acquired, but the estimate failed a declared validity or
  quality gate;
  \item \textbf{PREDICTED}: a dynamical model emitted an estimate without a new accepted
  observation; the horizon and source observation must remain attached;
  \item \textbf{MISSING}: no usable acquisition exists for that interval.
\end{itemize}

Conflating these states creates false continuity and biases downstream cognitive models.
Prediction can be valuable for control or display stabilization, but it is a model output, not a
retinal observation. A rejected aperture is also scientifically different from a missing interval:
the former contains photons whose registration was unreliable; the latter contains no usable
measurement.

\section{Current evidence and claim boundaries}

\subsection{Synthetic recovery}

The synthetic fixture begins with a single observed retinal frame smoothed by a 0.7-pixel Gaussian
to form a known-reference surrogate. It imposes known motion and calibrated synthetic noise, runs
the registration pipeline, and compares estimated and true displacement. This is the correct
environment for testing implementation logic and controlled failure modes. At the 24-line width,
6-line stride, DoG(1,8), spatial-NCC operating point, the current result is
\fixture{} 0.05221 arcmin RMSE with 393/393 good locks
(\texttt{synthetic\_precision\_24\_6}).

This is evidence that the pipeline can recover motion generated by the fixture. It is not evidence
that a human eye can be tracked at that numerical precision under the real device's optical blur, speckle,
scanner nonlinearity, calibration error, physiological deformation, and acquisition gaps. The
appropriate next step is to add realism one perturbation at a time and then test against independent
physical truth.

\begin{figure}[ht]
\centering
\includegraphics[width=\textwidth]{operating_curve.png}
\caption{Synthetic known-truth recovery. The fixture isolates algorithm behavior but does not
constitute biological or whole-device validation.}
\label{fig:synthetic}
\end{figure}

\subsection{Restricted 10 ms waveform fixture}

A second fixture injects minimum-jerk events into a real-frame-derived morphology at the measured
15.242 dB structure SNR. The noise model uses 15\% speckle power plus Gaussian noise and adds
two-dimensional Ornstein--Uhlenbeck drift. The event orientation is vertical along the
fast/image-row axis; it is not a two-axis microsaccade validation. At the 32-line width,
4-line stride operating point, the \fixture{} result recovers a 10 arcmin, 10 ms event in 3/3
trials under both 40.0 and 58.3 ms active-time brackets
(\texttt{synthetic\_10ms\_recovery\_32\_4}). Three trials are not enough to establish a reliable
recovery probability: the corresponding Wilson interval is 0.439--1.000.

\begin{figure}[ht]
\centering
\includegraphics[width=\textwidth]{waveform_recovery.png}
\caption{Restricted synthetic waveform recovery. The vertical 10 ms event, real-frame-derived
morphology, declared noise model, and three-trial count define the scope of the result.}
\label{fig:waveform}
\end{figure}

\subsection{Real retinal frame and pursuit fixture}

On the selected real retinal frame, 45 of 47 strips lock. This is a \measured{} lock result on that
capture and demonstrates that the algorithm finds a coherent trajectory through real texture.
No independent high-rate eye-motion truth is available for that frame, so biological accuracy and
precision cannot be calculated.

The longer pursuit demonstration compares SLO-derived motion with a lower-rate eye-camera record.
The \measured{} held-out correspondence is axis-specific:
$r_x=0.8729$, $r_y=0.8574$, with residual RMS of 16.73 and 23.73 arcmin, respectively
(\texttt{held\_out\_pursuit\_correlation}). The full result contains 66,058 active observations,
2,147 reacquisition observations, 17,391 predicted dead-time samples, 6,916 rejected mislocks, and
28,303 unlocked rows. These values are useful integration checks, but do not establish high-rate
SLO truth because the comparator has different bandwidth, calibration, latency, and spatial
precision. The 10 ms interval and its predicted samples come from an explicit unmeasured timing
hypothesis and must remain distinguishable from observed SLO samples.

\begin{figure}[ht]
\centering
\includegraphics[width=\textwidth]{real_demo.png}
\caption{Single-frame real-data demonstration. Accepted strips form a plausible trajectory, but
there is no independent high-rate biological truth for accuracy scoring.}
\label{fig:real}
\end{figure}

\begin{figure}[ht]
\centering
\includegraphics[width=\textwidth]{08_pursuit_result.png}
\caption{Longer pursuit demonstration and state timeline. Comparison with a lower-rate eye camera
is an integration check, not a high-rate truth standard.}
\label{fig:pursuit}
\end{figure}

\subsection{What remains unproved}

Current evidence does not yet establish:
\begin{itemize}
  \item repeatable whole-device precision on a physical model eye with commanded motion;
  \item traceable angular scale across subjects, eye position, field, and scanner direction;
  \item high-rate agreement with an independent reference over the relevant bandwidth;
  \item robustness to vessel-poor fields, blur, low signal, pathology, and large saccades;
  \item the measured scanner/blanking topology during the nominal inter-frame interval;
  \item patient exposure under continuous operation and single-fault conditions.
\end{itemize}

\section{Sensor requirements for cognitive signals}

\subsection{Four quantities that should not be collapsed}

The phrase ``high-speed eye tracking'' hides four independent questions.

\begin{enumerate}
  \item \textbf{Event rate}: how often events occur, such as microsaccades per second.
  \item \textbf{Waveform bandwidth}: how rapidly position or velocity changes within an event.
  \item \textbf{Acquisition rate}: how often the instrument samples sufficiently to estimate the
  desired timing or waveform with margin.
  \item \textbf{Spatial precision}: the smallest displacement that must be resolved after
  calibration and noise.
\end{enumerate}

A low event rate does not imply a low required acquisition rate. Microsaccades may occur only a few
times per second, yet millisecond-scale onset timing and direction estimation can require hundreds
of samples per second or more. Conversely, pupil dynamics may benefit from moderate frame rates but
do not require retinal arcsecond precision.

\subsection{Design envelopes}

\begingroup
\small
\setlength{\tabcolsep}{3pt}
\begin{longtable}{>{\bfseries}p{0.15\textwidth} p{0.18\textwidth} p{0.27\textwidth}
                  p{0.18\textwidth} p{0.13\textwidth}}
\caption{Signal-specific design envelopes. Rates are acquisition design ranges, not biological
event frequencies.}\label{tab:signals}\\
\toprule
Signal & Approximate scale & Acquisition interpretation & Precision regime & Likely sensor\\
\midrule
\endfirsthead
\toprule
Signal & Approximate scale & Acquisition interpretation & Precision regime & Likely sensor\\
\midrule
\endhead
Saccade & larger than 1--2$^\circ$ & 100--250 Hz detection; 200--500 Hz routine kinematics;
about 1000 Hz for fine onset/acceleration &
arcminute/coarse & camera or SLO\\
Microsaccade onset/rate & roughly 0.1--0.5$^\circ$ & at least 250 Hz; 500--1000 Hz for waveform
shape & coarse event detection & event/camera or SLO\\
Microsaccade direction & roughly 0.1--0.5$^\circ$ & approximately 500--1000 Hz when fine direction
and onset matter & arcminute, possibly finer at small amplitude & SLO or fusion\\
Drift & about 0--40 Hz content; reported velocity about 50--94 arcmin/s & ideal floor about 80 Hz;
200--500 Hz for credible velocity and spectra &
arcminute overall; short high-rate increments can become arcsecond-scale & SLO or fusion\\
Tremor & about 1 arcmin or smaller & 200 Hz is only the ideal Nyquist floor for 100 Hz content;
at least 500 Hz is a practical starting point; direct retinal work has used 1920 Hz &
sub-arcminute to arcsecond & SLO or specialized sensor\\
Pupil & millimeter-scale diameter & approximately 30--120 Hz for most cognitive pupillometry,
with protocol-specific filtering and latency & micrometer/coarse image scale & eye camera\\
\bottomrule
\end{longtable}
\endgroup

The table is a design synthesis from Rucci and Poletti \cite{RucciPoletti2015}, Rolfs
\cite{Rolfs2009}, Niehorster and colleagues \cite{Niehorster2020}, Poletti and colleagues
\cite{Poletti2015}, Bowers and colleagues \cite{Bowers2019}, and pupillography guidance
\cite{Kelbsch2019}; it is not a universal specification. Sampling targets depend on the
analysis objective, latency budget, estimator, noise, anti-alias filtering, and desired uncertainty.
For tremor, 200 Hz is the mathematical Nyquist floor for a 100 Hz component, not a robust
instrument requirement. Bowers and colleagues sampled retinal motion at 1920 Hz, enabling analysis
through 960 Hz \cite{Bowers2019}; that acquisition choice is evidence of experimental practice, not
a requirement that every cognitive mode operate at 1920 Hz.

\begin{figure}[ht]
\centering
\includegraphics[width=\textwidth]{07_signal_landscape.png}
\caption{Temporal acquisition demand and spatial precision are orthogonal. Tremor is intrinsically
arcsecond-scale; fine microsaccade direction and short high-rate drift increments can also become
precision-limited.}
\label{fig:signals}
\end{figure}

\subsection{Implication for sensor fusion}

No single sensor must dominate all modes. A pupil camera provides robust coarse gaze and pupil
diameter over a wide field. The SLO provides retinal texture and potentially higher relative spatial
precision during valid acquisition. OPM-MEG, stimulus, and behavior provide additional synchronized
modalities. A production platform should preserve hardware time and uncertainty so models can use
the best available sensor rather than forcing every signal through one image stream.

This corrects an overly strong slogan that ``only tremor needs the SLO.'' Tremor is the clearest
intrinsically arcsecond-scale signal, but fine direction estimates for small microsaccades and
short-interval drift increments can also demand very small displacement precision. Whether this
specific SD-SLO reaches those regimes remains an empirical question.

\section{Inter-frame dead time and continuous acquisition}

\subsection{What is known and what is not}

Two legitimate but incompatible conditional timing models must be kept separate.

\begin{enumerate}
  \item \textbf{Legacy capture bracket.} The pursuit result explicitly assumes a 68.3387 ms frame
  period split into 58.3387 ms candidate active time and 10.0 ms dead time. The conditional dead
  fraction is $10/68.3387=14.63\%$. This is a \hypothesis{}, not scanner telemetry.
  \item \textbf{Candidate-clock cross-design calculation.} Applying the technical note's
  20 kline/s bidirectional clock to the legacy 808-line geometry gives 40.4 ms active-line time.
  At the measured 14.633 fps frame period, 27.9387 ms remains unexplained, a conditional fraction
  of 40.88\%. This mixes a candidate fast-axis clock with legacy frame geometry and therefore is
  a \derived{} diagnostic, not a measured duty cycle.
\end{enumerate}

The repository does not contain measured slow-axis position, illumination state, detector validity,
or complete scanner telemetry for either interval. Calling either interval ``flyback'' would imply
a cause that has not been demonstrated. The correct term is \emph{uncharacterized inter-frame
dead time}.

\begin{figure}[ht]
\centering
\includegraphics[width=\textwidth]{06_flyback_gap.png}
\caption{Conditional blind interval from the candidate 20 kline/s clock combined with legacy
808-line geometry and frame rate. This cross-design calculation does not identify its physical
cause; scanner and blanking telemetry are required.}
\label{fig:deadtime}
\end{figure}

For event onsets uniform in phase relative to the frame cycle, the fraction beginning in an
unobserved interval equals that interval's conditional fraction: 14.63\% for the explicit legacy
bracket or 40.88\% for the candidate-clock cross-design calculation. These are conditional
consequences of separate models, not measurements of the present instrument.

\subsection{Why prediction does not remove the blind interval}

A dynamical model can propagate the last accepted state into a gap:
\[
  \hat{\mathbf x}_{k|k-1}=\mathbf F_k\hat{\mathbf x}_{k-1|k-1},\qquad
  \mathbf P_{k|k-1}=\mathbf F_k\mathbf P_{k-1|k-1}\mathbf F_k^\mathsf T+\mathbf Q_k.
\]
Uncertainty grows because no new retinal innovation constrains the state. A saccade or
microsaccade that begins inside the gap can violate the motion model. Prediction is therefore useful
for control continuity, but it cannot establish event onset or direction as an observation.

\subsection{Freeze the line-record interface before choosing a scanner topology}

The no-regret architectural move is to define a timestamped line record that remains valid across
scanner options:
\[
  \mathcal L_i =
  \{t_i,\mathrm{line\ id},\mathrm{direction},\mathbf p_i,
    \mathrm{illumination},\mathrm{validity},\mathrm{calibration},\mathrm{revision}\}.
\]
Here $\mathbf p_i$ can be a measured position trace or reference to a versioned LUT. The interface
does not promise gapless observation and does not presuppose resonant, sawtooth, dual-axis, or
fused-sensor hardware. It simply prevents scanner topology from being hidden inside an opaque frame
array.

Candidate options can then be compared against one interface: characterize and reduce the present
dead time; acquire useful lines through a return path if optically and safely valid; use a scanner
topology designed for higher duty cycle; or fuse SLO with a continuously observing coarse sensor.
Selection should be based on measured valid-line duty cycle, timestamp error, spatial calibration,
registration performance, safety margin, OPM coexistence, mechanics, cost, and production risk.

\section{Laser-safety framework for continuous operation}

\subsection{Three evaluations that must remain distinct}

``Laser safety'' is not one calculation.

\begin{enumerate}
  \item \textbf{Laser-product classification.} IEC 60825-1 and applicable U.S. 21 CFR provisions
  address accessible emission, product class, labeling, engineering controls, and scanning
  safeguards \cite{IEC60825,FDA1040}.
  \item \textbf{Ophthalmic patient exposure.} ISO 15004-2 and FDA-recognized ANSI Z80.36 address
  optical-radiation hazards for ophthalmic instruments under normal and relevant fault
  conditions \cite{ISO15004,ANSI8036}.
  \item \textbf{Medical-device classification and pathway.} FDA product codes and recognized
  standards govern the regulatory pathway. This classification is not equivalent to an IEC laser
  class \cite{FDAMYC}.
\end{enumerate}

The applicable clauses, wavelength weighting, angular subtense treatment, scanning correction,
measurement aperture, and averaging windows must be selected by qualified laser-safety and
regulatory personnel for the actual product configuration. This paper is an engineering input and
not a safety determination.

\subsection{Parameterized exposure quantities}

The fundamental quantities include radiant power $P(\lambda,t)$ at the cornea, retinal irradiance,
radiant exposure, dwell, revisit pattern, and cumulative session history. In simplified notation,
\[
  E_{\mathrm{ret}}(\lambda,t)=
  \frac{P(\lambda,t)\,T_{\mathrm{eye}}(\lambda)\,C_{\mathrm{geom}}(t)}
       {A_{\mathrm{ret}}(\lambda,t)},
\qquad
  H_{\mathrm{ret}}=\int E_{\mathrm{ret}}(t)\,dt .
\]
This form is intentionally generic. A standards calculation may use weighted source radiance or
radiance dose, apparent-source angular subtense, wavelength-dependent hazard functions, pulse-group
rules, and scanning-failure provisions rather than this simple area model.

For one fast sweep with retinal width $W$ and center velocity $v_{\max}$, a crude center dwell over
spot diameter $d$ is $t_d\approx d/v_{\max}$. Near sinusoidal turnarounds, velocity approaches zero,
so the center approximation is invalid. Real exposure must integrate the measured two-axis
trajectory, beam profile, blanking state, and fault behavior.

\subsection{Inputs that are still required}

A defensible calculation requires, at minimum:
\begin{itemize}
  \item measured spectral power at the corneal plane for every operating mode;
  \item wavelength distribution and source linewidth;
  \item beam diameter, divergence, and profile at the standard measurement location;
  \item complete measured fast- and slow-axis position-versus-time traces;
  \item illumination/blanking waveform and detector-validity timing;
  \item nominal session duration, repetition schedule, and cumulative-use model;
  \item scanner tolerances and credible single-fault trajectories;
  \item evidence that independent hardware removes or limits emission under scanner faults;
  \item identification of the exact standard editions, clauses, corrections, and regulatory
  pathway used by the reviewer.
\end{itemize}

Those values are presently \missing{} in the evidence ledger. Consequently, no allowed-power
number or continuous-operation safety conclusion is stated here.

\begin{figure}[ht]
\centering
\includegraphics[width=\textwidth]{09_safety_inputs.png}
\caption{Safety analysis is a parameterized measurement and fault-control problem. Geometry alone
cannot establish compliance.}
\label{fig:safety}
\end{figure}

\subsection{Controls should not depend on application software}

Candidate controls include independent scanner-motion monitoring, hardwired emission shutdown,
power limiting, startup self-test, watchdogs, blanking verification, session-dose accounting,
fault logging, and conservative stop states. Their response time and independence must be derived
from the applicable hazard analysis. If scanner position, clock validity, or blanking state becomes
invalid, the safe action must not rely on a high-level gaze process continuing to execute.

\section{Production architecture for a cognitive platform}

\subsection{Keep, change, configure, and defer}

\begingroup
\setlength{\tabcolsep}{4pt}
\begin{longtable}{>{\bfseries}p{0.15\textwidth} p{0.34\textwidth} p{0.38\textwidth}}
\caption{Disposition of major legacy assumptions.}\label{tab:ledger}\\
\toprule
Disposition & Element & Reason\\
\midrule
\endfirsthead
\toprule
Disposition & Element & Reason\\
\midrule
\endhead
Keep & Confocal retinal imaging, optical calibration, reference registration, safety engineering,
portable mechanics, controlled OPM coexistence & These capabilities support both retinal health
and cognitive measurements.\\
Change & Opaque frame-only acquisition, uncharacterized dead time, hidden timestamp assumptions,
and state interpolation that erases provenance & They prevent event-level interpretation and
traceable multimodal alignment.\\
Configure & Strip width/stride, optical power mode, field, sensor fusion, reference policy,
display protocol, and estimator latency & Different experiments need different time/space/safety
trade-offs; one hard-coded operating point is not optimal.\\
Defer & Final scanner topology, universal arcsecond claim, long-session optical mode, and a single
ML representation & Current evidence does not justify freezing these decisions.\\
\bottomrule
\end{longtable}
\endgroup

\subsection{Layered data and timing model}

A production cognitive instrument should expose:
\begin{enumerate}
  \item \textbf{hardware time}: one monotonic clock or a calibrated mapping among clocks;
  \item \textbf{sample records}: detector samples, scanner telemetry, illumination, and validity;
  \item \textbf{line records}: direction, time interval, position/LUT version, and calibration;
  \item \textbf{observation records}: displacement, quality, aperture, and reference identity;
  \item \textbf{estimate records}: filtering/prediction state and uncertainty;
  \item \textbf{multimodal records}: display photons or photodiode, eye camera, behavior, OPM-MEG,
  and experiment markers;
  \item \textbf{provenance}: hardware, firmware, software, calibration, safety mode, and model
  revision identifiers.
\end{enumerate}

Frames can still be generated for clinical review and visualization, but they become one product
of the line stream rather than the only representation. This preserves compatibility with the
retinal-health use case while making time explicit for cognition.

\subsection{Synchronization requirement}

Each subsystem clock $t_j$ should be mapped into a master timebase:
\[
  t_{\mathrm{master}} = a_j t_j + b_j + \epsilon_j(t),
\]
where $a_j$ captures rate error, $b_j$ captures offset, and $\epsilon_j$ captures residual jitter or
nonlinearity. Synchronization acceptance must report both mean error and tails across the full
session. A software timestamp taken after buffering is insufficient evidence of photon or stimulus
time. Display timing should be measured with a photodiode or equivalent optical marker when the
experiment depends on actual photon onset.

\section{Validation and production freeze gates}

\subsection{Validation sequence}

The validation program should advance in causal order.

\begin{enumerate}
  \item \textbf{Timing topology.} Record fast/slow position, detector clock, illumination, blanking,
  frame markers, and validity on common hardware time. Explain every interval.
  \item \textbf{Optical/electronic transfer.} Measure spectrum, power, spot/beam profile, field,
  detector response, ADC noise, delay, and forward/reverse differences.
  \item \textbf{Model-eye motion.} Command traceable position trajectories spanning drift,
  microsaccade-like steps, saccades, and tremor-band components.
  \item \textbf{Registration operating curves.} Sweep strip width, stride, signal, blur, vessel
  content, search range, and motion; report error, lock, false lock, latency, and uncertainty.
  \item \textbf{Biological comparison.} Compare with an independent high-rate reference whose own
  bandwidth, latency, calibration, and uncertainty are characterized.
  \item \textbf{Multimodal synchronization.} Measure SLO, camera, display, behavior, and OPM
  alignment end to end, including drift over session duration.
  \item \textbf{Safety and fault verification.} Complete standards analysis and test scanner,
  blanking, watchdog, power, and shutdown faults.
  \item \textbf{Production repeatability.} Repeat across units, operators, subjects, sessions, and
  controlled revisions.
\end{enumerate}

\subsection{Candidate measurable gates}

\begin{table}[ht]
\centering
\caption{Production freeze gates. Numerical thresholds should be approved by the responsible
engineering and scientific owners before they are treated as specifications.}
\begin{tabularx}{\textwidth}{>{\bfseries}p{0.18\textwidth} X}
\toprule
Gate & Required evidence\\
\midrule
G0: use cases & Named scientific observables, uncertainty tolerance, session length, and regulatory
mode for each operating configuration.\\
G1: timing & Measured line/frame topology, valid duty cycle, timestamp error distribution, and
explained dead intervals.\\
G2: optics/electronics & Measured power/spectrum, field and spot profile, transfer function, delay,
noise, and direction asymmetry.\\
G3: calibration & Traceable spatial/time transform with versioning and repeatability across field,
direction, and operating modes.\\
G4: gaze & Model-eye and biological error, precision, lock, false-lock, latency, and state
classification by motion/quality regime.\\
G5: safety & Qualified review, normal/single-fault calculations, independent controls, and verified
fault responses for every enabled optical mode.\\
G6: OPM & Controlled coexistence record and revision matrix tied to the production configuration.\\
G7: multimodal & End-to-end display, SLO, camera, behavior, and OPM synchronization with uncertainty
and session drift.\\
G8: data/ML & Stable schemas, immutable raw records, calibration/revision provenance, leakage-safe
splits, and external validation.\\
\bottomrule
\end{tabularx}
\end{table}

\subsection{A 30/60/90-day decision sequence}

In the first 30 days, instrument timing and optical state should be measured on a common clock, the
line-record schema frozen at version 0.1, and all current captures made ingestible without
discarding invalid intervals. In parallel, responsible owners should define cognitive use cases,
session durations, and safety operating modes.

By 60 days, a traceable model-eye motion rig should exercise the full registration chain. Candidate
scanner and fusion options should be scored with the same metrics and interface. Preliminary
normal/fault optical measurements should identify whether a continuous mode is plausible and what
controls it requires.

By 90 days, the program should be able to make evidence-backed keep/change/defer decisions on
scanner topology, optical duty cycle, sensor fusion, calibration hardware, timing architecture,
and production data contracts. Decisions that still lack measured evidence should remain open
rather than being converted into assumptions by schedule pressure.

\section{Limitations}

This analysis has several explicit limitations.
\begin{itemize}
  \item The technical-note parameters describe a candidate scanner and sampling design; they are
  not all measurements of the current instrument.
  \item Real-data demonstrations lack independent high-rate biological truth.
  \item Synthetic performance depends on the fixture's motion, texture, noise, blur, and
  calibration assumptions.
  \item Nominal retinal scale does not replace subject-specific biometry or measured angular
  calibration.
  \item Literature from adaptive-optics or specialized retinal systems does not establish the
  precision of this SD-SLO.
  \item The cause and optical state of the inter-frame interval are unmeasured in the available
  evidence.
  \item The safety section identifies standards regimes, equations, inputs, and controls but is
  not a compliance conclusion or allowed-power calculation.
  \item OPM coexistence is accepted as a program fact; production applicability still requires a
  controlled revision record.
\end{itemize}

\section{Conclusion}

An SLO can provide more than retinal pictures because each detector sample is tied to scanner
position and time. Eye motion deforms the raster; short-strip registration can invert that
deformation into a time-resolved relative trajectory. The candidate 10 kHz sinusoid supports a
20 kline/s clock, and the 24/6 strip example produces 1.20 ms apertures every 0.30 ms during active
acquisition. These outputs are overlapping and correlated, so output rate must not be confused with
independent information bandwidth.

The current software evidence is promising but bounded. Synthetic fixtures validate algorithm
behavior under known assumptions. Real retinal captures show that vessel-based registration can
lock and track coherent motion. Neither establishes whole-device arcsecond performance or
high-rate biological accuracy. Those claims require traceable model-eye motion, independent
high-rate reference measurements, calibrated transforms, and explicit state/provenance records.

The program's pivot from intermittent retinal imaging to a cognitive platform makes time and
missingness first-class design variables. The most durable production decision is to freeze a
timestamped line-record and multimodal timing interface while leaving scanner topology and optical
duty cycle open until they are measured. Frames remain useful outputs, but should no longer hide
the temporal structure of acquisition.

Finally, continuous illumination cannot be justified by a duty-cycle argument. Laser-product
classification, patient exposure, and medical-device pathway are distinct evaluations. The
necessary measured spectrum, power, beam, trajectory, session, and fault inputs must be completed
and reviewed before a safety conclusion. With those boundaries preserved, the existing retinal
capability can become the precision component of a broader synchronized cognitive instrument
without overclaiming what has already been demonstrated.

\appendix
\section{Reproducibility and source artifacts}

The primary engineering source is
\textit{SLO Scan Bandwidth Sampling Technical Note v2}, identified in the evidence ledger by
SHA-256:
\begin{center}
  \small\texttt{e9fb9728c77d352bcd55ccb6eb66979e}\\
  \small\texttt{e415b1c9d416d32b517f1c924a6c33c4}
\end{center}
Key repository artifacts are:
\begin{itemize}
  \item \path{cognitive_gaze_evidence.json}: claim and input ledger;
  \item \path{cognitive_gaze_figure_manifest.json}: figure/input hashes;
  \item \path{build_cognitive_gaze_figures.py}: figure generator;
  \item \path{build_cognitive_gaze_deck.py}: PowerPoint generator;
  \item \path{src/intraframe_ncc/}: registration and trajectory implementation;
  \item \path{src/sdslo_gaze/}: scan timing, desinusoiding, dynamics, contracts, and evaluation.
\end{itemize}

\section{Recommended minimum line-record fields}

\begingroup
\setlength{\tabcolsep}{4pt}
\begin{longtable}{>{\ttfamily}p{0.24\textwidth} p{0.65\textwidth}}
\toprule
\normalfont\bfseries Field & \bfseries Meaning\\
\midrule
timestamp\_start/end & Hardware-clock interval covered by the line.\\
line\_id / frame\_id & Monotonic line identity; optional frame grouping.\\
direction & Forward/reverse or other scanner phase identity.\\
position\_trace & Measured scanner positions or versioned LUT reference.\\
illumination\_state & On/off/power mode and relevant blanking state.\\
detector\_validity & Saturation, clipping, packet loss, or analog/ADC faults.\\
calibration\_id & Spatial, delay, direction, and distortion transform version.\\
hardware\_revision & Optics, scanners, detector, electronics, and mechanics.\\
firmware/software & Acquisition and processing versions.\\
safety\_mode & Approved optical configuration and control state.\\
\bottomrule
\end{longtable}
\endgroup

\begin{thebibliography}{99}

\bibitem{Bennett1994}
A. G. Bennett, M. Rudnicka, and D. F. Edgar,
``Improvements on Littmann's method of determining the size of retinal features by fundus
photography,'' \emph{Graefe's Archive for Clinical and Experimental Ophthalmology},
vol. 232, pp. 361--367, 1994.

\bibitem{Bowers2019}
N. R. Bowers et al.,
``The effects of fixational tremor on the retinal image,''
\emph{Journal of Vision}, vol. 19, no. 11, article 8, 2019,
doi:10.1167/19.11.8.

\bibitem{Kelbsch2019}
C. Kelbsch et al.,
``Standards in pupillography,''
\emph{Frontiers in Neurology}, vol. 10, 2019.

\bibitem{RucciPoletti2015}
M. Rucci and M. Poletti,
``Control and functions of fixational eye movements,''
\emph{Annual Review of Vision Science}, vol. 1, pp. 499--518, 2015,
doi:10.1146/annurev-vision-082114-035742.

\bibitem{Poletti2015}
M. Poletti, J. Aytekin, and M. Rucci,
``Head-eye coordination at a microscopic scale,''
\emph{Current Biology}, 2015,
doi:10.1016/j.cub.2015.11.004.

\bibitem{Rolfs2009}
M. Rolfs,
``Microsaccades: small steps on a long way,''
\emph{Vision Research}, vol. 49, no. 20, pp. 2415--2441, 2009.

\bibitem{Niehorster2020}
D. C. Niehorster et al.,
``The Tobii Pro Spectrum: a useful tool for studying microsaccades?''
\emph{Behavior Research Methods}, 2020,
doi:10.3758/s13428-020-01430-3.

\bibitem{IEC60825}
International Electrotechnical Commission,
\emph{IEC 60825-1: Safety of Laser Products---Part 1: Equipment Classification and Requirements},
applicable edition.

\bibitem{ISO15004}
International Organization for Standardization,
\emph{ISO 15004-2:2024 Ophthalmic Instruments---Fundamental Requirements and Test Methods---Part 2:
Light Hazard Protection}.

\bibitem{ANSI8036}
American National Standards Institute,
\emph{ANSI Z80.36-2021: Ophthalmics---Light Hazard Protection for Ophthalmic Instruments}.

\bibitem{FDA1040}
U.S. Food and Drug Administration,
\emph{21 CFR 1040.10 and 1040.11; Laser Notice No. 56}, applicable requirements.

\bibitem{FDAMYC}
U.S. Food and Drug Administration,
\emph{Product Classification: Scanning Laser Ophthalmoscope, Product Code MYC}.

\end{thebibliography}

\end{document}
